\section{第一条扩展路径：预训练 Scaling Laws}

\subsection{三个共同扩展的轴}

预训练阶段的经验规律是：在合理的数据与优化配置下，增加\textbf{训练计算量、训练数据规模和模型参数量}通常都会降低测试损失。常用的抽象形式是幂律：

$$
L(x)=L_{\infty}+A x^{-\alpha},\qquad \alpha>0.
$$

\begin{itemize}
    \item $L(x)$：给定资源规模 $x$ 时的测试损失；
    \item $L_{\infty}$：该数据分布和建模假设下难以消除的损失下界；
    \item $A$：与任务、数据和模型族相关的比例常数；
    \item $x$：可以表示计算量、数据量或参数量；
    \item $\alpha$：幂律指数，决定资源增加后损失下降的速度。
\end{itemize}

幂律的重要性不在于“永远有效”，而在于它曾经让团队能够在训练前预测更大模型的收益，并据此投入更大规模的算力与数据工程。

\begin{figure}[H]
\centering
\includegraphics[width=0.96\textwidth]{figures/scaling-laws.jpg}
\caption{预训练 scaling laws：计算、数据和参数三个扩展轴。\protect\footnotemark}
\end{figure}
\footnotetext{视频讲解区间：00:02:50--00:03:53。}

\begin{knowledgebox}{“更大”不是只指参数更多}
一个参数极大的模型如果数据不足，会重复记忆有限样本；如果训练计算不足，参数得不到充分优化；如果数据质量差，额外计算只会更充分地拟合噪声。因此现代 scaling 的核心是\textbf{参数、token 和有效计算的协同配置}。
\end{knowledgebox}

\subsection{从 BERT、GPT-2 到 GPT-4：规模成为能力来源}

2018--2024 年间，主流语言模型的规模快速增长。课程用这段历史强调两点：第一，扩大规模带来了连续的平均性能提升；第二，某些能力并非在小模型上平滑可见，而是在达到一定规模后突然成为可观测行为。

\begin{figure}[H]
\centering
\includegraphics[width=0.92\textwidth]{figures/model-growth.jpg}
\caption{大语言模型参数规模的快速增长。\protect\footnotemark}
\end{figure}
\footnotetext{视频讲解区间：00:03:55--00:04:46。}

这里应避免一个常见误解：图中参数规模增长并不意味着参数是唯一因果变量。数据配比、架构、优化器、训练稳定性和后训练方法都在同期变化，因此历史曲线表达的是\textbf{系统工程整体扩展的结果}。

\subsection{Zero-shot、Few-shot 与上下文学习}

GPT-3 之后，一个关键变化是模型可以直接从 prompt 中理解任务：

\begin{itemize}
    \item \textbf{Zero-shot}：只给任务说明和输入，不给示例；
    \item \textbf{One-shot}：额外给一个输入--输出示例；
    \item \textbf{Few-shot}：给少量示例，让模型从上下文推断任务映射。
\end{itemize}

这类能力被称为 \textbf{in-context learning}：模型权重没有更新，但上下文中的模式改变了本次前向计算的行为。

\begin{figure}[H]
\centering
\includegraphics[width=0.94\textwidth]{figures/zero-few-shot.jpg}
\caption{Zero-shot 与 few-shot 提示的结构差异。\protect\footnotemark}
\end{figure}
\footnotetext{视频讲解区间：00:05:40--00:06:38。}

\begin{warningbox}{上下文学习不等于参数学习}
Few-shot 示例不会在推理时永久写入权重。它更像为当前请求建立一个临时任务解释。上下文结束后，这个“学习结果”通常不会跨请求保留，除非系统另有外部记忆或继续训练机制。
\end{warningbox}

\subsection{Chain-of-Thought：把中间过程也放进上下文}

Chain-of-thought（CoT）提示的关键不是多写几句话，而是提供\textbf{从问题到答案的中间推理轨迹}。相比只展示最终答案，模型可以模仿问题分解、状态更新和局部计算方式。

\begin{figure}[H]
\centering
\includegraphics[width=0.96\textwidth]{figures/chain-of-thought.jpg}
\caption{标准提示与 chain-of-thought 提示的对比。\protect\footnotemark}
\end{figure}
\footnotetext{视频讲解区间：00:07:18--00:09:01。}

课程强调，CoT 效果随模型规模增长而明显增强。在小模型上，给出推理示例可能只增加模仿负担；在足够大的模型上，它能触发更可靠的多步计算。

\begin{figure}[H]
\centering
\includegraphics[width=0.92\textwidth]{figures/cot-emergence.jpg}
\caption{Chain-of-thought 能力在更大模型上出现明显跃迁。\protect\footnotemark}
\end{figure}
\footnotetext{视频讲解区间：00:09:01--00:10:18。}

\begin{importantbox}{从“预测下一个 token”到“执行一个计算过程”}
CoT 表明，同一个 next-token predictor 可以在上下文中表现得像一个迭代计算器。模型并未改变目标函数，但通过生成中间 token，把一次困难映射拆成多个更容易的局部映射。后续的 reasoning model 与 Agent 都沿用了这个思想：\textbf{允许系统在给出最终答案前消耗更多中间步骤。}
\end{importantbox}

\subsection{本章小结}

预训练 scaling 同时扩大了模型的平均能力和可被 prompt 激活的潜在行为。Zero-shot、few-shot 与 CoT 展示了上下文如何在不更新权重的情况下重新组织模型计算，但模型“知道什么”与“按人类意图行动”仍是两件事。

